\section{Криволинейные интегралы}
\subsection{Кривые в $\R^n$}
\begin{definition}\tab
    \begin{enumerate}
        \item Непрерывное отображение $\phi: [a,b]\to \R^k$ такое, что $\forall t\in [a,b]:\\ \phi(t)=(x_1(t), \dots, x_k(t))$ называется параметризацией кривой.
        \item Пусть $\phi_1:[\alpha_1, \beta_1]\to \R^k,\ \phi_2:[\alpha_2, \beta_2]\to \R^k$. Считаем, что $\phi_1$ и $\phi_2$ эквивалентны (параметризуют одну и ту же кривую), если существует строго монотонная функция $g: [\alpha_1, \beta_1]\to [\alpha_2, \beta_2]$ такая, что $\phi_2\equiv \phi_1\circ g$ на $[\alpha_1, \beta_1]$.
        \item Кривой называется класс эквивалентности параметризаций по указанному отношению.
    \end{enumerate}
\end{definition}
\begin{definition}\tab
    \begin{enumerate}
        \item Кривая называется простой, если всякая ее параметризация инъективна.
        \item Кривая называется замкнутой, если для любой ее параметризации \\
        $\phi: [\alpha,\beta]\to \R^k$ выполнено равенство $\phi(\alpha)=\phi(\beta)$.
        \item Кривая называется простой замкнутой, если из равенства $\phi(t_1)=\phi(t_2)$ следует, что либо $t_1= t_2$, либо $\{t_1, t_2\}=\{\alpha,\beta\}$.
    \end{enumerate}
\end{definition}
\begin{definition}
    Пусть $T=\{t_i\}$ --- разбиение отрезка $[\alpha, \beta]$:
    \[\alpha=t_0<t_1<t_2<\dots<t_n=\beta\]
    Тогда множество
    \[\Lambda(T)=\bigcup\limits_{i=1}^n [\phi(t_{i-1}), \phi(t_i)]\]
    называется вписанной ломаной в кривую $\gamma$.
\end{definition}
\begin{definition}
    Число
    \[|\Lambda(T)|=\sum\limits_{i=1}^{n} |\phi(t_i)-\phi(t_{i-1})|\]
    называется длиной вписанной ломаной, а число
    \[|\gamma|=\sup\limits_T |\Lambda(T)|\]
    называется длиной кривой. Если длина кривой конечна, то кривая называется спрямляемой.
\end{definition}
\begin{reminder}
    Функция $f$, определенная на промежутке $I$ имеет ограниченную вариацию, если
    \[\var(f,I)=\sup\limits_{n\in \N} \sum\limits_{i=1}^{n}|f(t_i)-f(t_{i-1})|<+\infty\]
\end{reminder}
\begin{theorem} [Об условиях спрямляемости кривых] $\\$
    Следующие условия эквивалентны:
    \begin{enumerate}
        \item $|\gamma|<+\infty$.
        \item Всякая параметризация кривой $\gamma$ имеет ограниченную вариацию.
        \item Существует параметризация кривой $\gamma$ с ограниченой вариацией.
    \end{enumerate}
\end{theorem}
\begin{proof}\tab
    \begin{enumerate}
        \item $(1) \Rightarrow (2)$:
        Пусть $\phi:[\alpha,\beta]\to \R^k$ --- произвольная параметризация\\
        кривой $\gamma$, $T$ --- разбиение отрезка $[\alpha,\beta]$. Тогда $|\Lambda(T)|<|\gamma|=\text{const}$
        \begin{multline*}
            |\gamma|\geq\sum\limits_{i=1}^{n}|\phi(t_i)-\phi(t_{i-1})|=\\
            =\sum\limits_{i=1}^{n}\sqrt{(x_1(t_i)-x_1(t_{i-1}))^2+\dots+(x_k(t_i)-x_k(t_{i-1}))^2}\geq \\
            \geq \sum\limits_{i=1}^{n}|x_j(t_i)-x_j(t_{i-1})|
        \end{multline*}
        \item $(2) \Rightarrow (3)$: Очевидно.
        \item $(3) \Rightarrow (1)$: Пусть существует параметризация $\phi:[\alpha,\beta]\to \R^k$ и ее вариация ограничена. Тогда
        \begin{multline*}
            |\Lambda(T)|=\sum\limits_{i=1}^{n}|\phi(t_i)-\phi(t_{i-1})|\leq\\
            \sum\limits_{i=1}^{n}\sum\limits_{j=1}^{k}|x_j(t_i)-x_j(t_{i-1})|=\sum\limits_{j=1}^{k}\sum\limits_{i=1}^{n}|x_j(t_i)-x_j(t_{i-1})|\leq\\
            \leq \sum\limits_{j=1}^{k}\text{Var}(x_j([\alpha,\beta]))
        \end{multline*}
    \end{enumerate}
\end{proof}
\begin{example} [Неспрямляемой кривой] $\\$
    Построим непрерывную функцию $\gamma=(x(t),y(t))$ неограниченной вариации:
    \[x(t)\equiv t,\ \ y(t)=\begin{cases}
        0,\ t=0,\\
        t\cdot \sin{(\frac{\pi}{t})},\ t\in [0,1]
    \end{cases}\]
    %тут будет картинка.
\end{example}
\begin{exercise}
    Приведите пример простой кривой (без самопересечений) положительной меры Жордана.
\end{exercise}
\begin{reminder}
    Пусть у кривой $\gamma$ существует параметризация $\phi\in \mathcal{C}^1$. Тогда
    \[|\gamma|=\int\limits_{\alpha}^{\beta} \sqrt{(x_1'(t))^2+\dots+(x_k'(t))^2}\ dt\]
\end{reminder}
\begin{comm}
    Отображения из класса $\mathcal{C}^1$, у которых вектор скорости не обращается в ноль, в этом курсе будем называть гладкими.
\end{comm}
\subsection{Криволинейный интеграл первого рода}
\begin{reminder}\tab
    \begin{enumerate}
        \item Пусть $f$ и $g$ определены на $[a,b], T=\{t_i\}$ и $H=\{\xi_i\}$ --- разбиение и разметка отрезка $[a,b]$ соответственно. Тогда интегральной суммой Римана-Стилтьеса назовем
        \[\sigma_s(f,T,H)=\sum\limits_{i=1}^{n}f(\xi_i)\cdot (g(t_i)-g(t_{i-1}))\]
        \item Если 
        \[\forall \epsilon>0\ \exists\ \delta>0,\ d(T)<\delta: |\sigma_s(f,T,H)-I|<\epsilon\]
        то
        \[I=\int\limits_{a}^{b}f(x)\ d(g(x))\]
        называется интегралом Римана-Стилтьеса.
    \end{enumerate}
\end{reminder}
\begin{definition}
    Пусть $f$ определена на $\phi([\alpha,\beta])$, а кривая $\gamma$ спрямляема. Криволинейным интегралом первого рода назовем следующий интеграл Римана-Стилтьеса:
    \[\int\limits_{\gamma}f(x_1,\dots,x_k)\ ds=\int\limits_{\alpha}^{\beta}f(x_1(t),\dots,x_k(t))\ ds(t)\]
    где $s(t_0)$ --- длина образа $[\alpha,t_0],\ t_0\in [\alpha,\beta]$.
\end{definition}
\begin{definition}[Альтернативное определение криволинейного интеграла первого рода]\tab
    \begin{enumerate}
        \item Пусть $T=\{t_i\}$ --- разбиение отрезка $[\alpha,\beta],\ \Delta_i=[\phi(t_{i-1}),\phi(t_i)]$. Множество $H=\{\xi_i\}$, где $\xi_i$ выбран на дуге, которую стягивает отрезок $\Delta_i$, называется разметкой кривой. Диаметром вписанной кривой назовем $d(\Lambda(T))$ --- наибольшую из длин отрезков в составе $\Lambda(T)$.
        \item Выражение
        \[\sigma(f,\Lambda(T), H)=\sum\limits_{i=1}^{n}f(\xi_i)\cdot |\Delta_i|\]
        назовем интегральной суммой для криволинейного интеграла.
        \item Наконец, если
        \[\forall \epsilon>0\ \exists\ \delta>0,\ d(\Lambda(T))<\delta: |\sigma(f,\Lambda(T), H)-I|<\epsilon\]
        то 
        \[I=\int\limits_{\gamma}f(x_1,\dots,x_k)\ ds\]
    называется криводинейным интегралом первого рода.
    \end{enumerate}
\end{definition}
\begin{lemma}
    Если для функции $f$ выполнено одно из определений криволинейного интеграла первого рода для простых кривых, то $f$ ограничена на $\phi([\alpha,\beta])$.
\end{lemma}
\begin{proof}
    \[\sigma_{s}(f,T,\eta)=\sum\limits_{i=1}^{n}f(\phi(\eta_i))(s(t_i)-s(t_{i-1}))\]
\[\forall \epsilon>0\ \exists\ \delta>0,\ d(T)<\delta: |\sigma_s-I|<\epsilon\]
\[I-\epsilon<\sigma_s<I+\epsilon\]
Пусть
\[c=\sum\limits_{i\neq m}f(\phi(\eta_i))(s(t_i)-s(t_{i-1}))\]
Итак
\[I-c-\epsilon<f(\phi(\eta_m))(s(t_m)-s(t_{m-1}))<I-c+\epsilon\]
\end{proof}
\begin{lemma}
    Пусть $\gamma$ --- спрямляемая кривая, $\phi$ --- ее параметризация. Тогда 
    \[\forall \epsilon>0\ \exists\ \delta>0,\ d(\Lambda(T))<\delta: |\gamma|-|\Lambda(T)|<\epsilon\]
\end{lemma}
\begin{proof} Без доказательства.
    % $\exists\ T': |\gamma|-|\Lambda(T')|<\frac{\epsilon}{2}$.
    % Поймем насколько меняется длина ломаной, если добавить в разбиение отрезка еще одну точку $\tau$. Для $\tau_1,\tau_2\in [\alpha,\beta]$ обозначим
    % \[\omega(\delta)=\sup\limits_{|tau_1-\tau_2|<\delta}\ \|\phi(\tau_1)-\phi(\tau_2)\|\]
    % и заметим, что $\omega(\delta)\to 0,\ \delta\to 0+ \Leftrightarrow \phi$ --- раномерно непрерывна.
    % \[\forall \epsilon>0\ \exists\ \delta>0,\ |\tau_1-\tau_2|<\delta: \|\phi(\tau_1)-\phi(\tau_2)\|<\theta=\frac{\epsilon}{4N}\]
\end{proof}
\begin{theorem}
    Два определение криволинейных интегралов первого рода эквивалентны.
\end{theorem}
\begin{proof} \textit{(не спрашивают на экзамене)}\tab
    \begin{enumerate}
        \item $(1)\Rightarrow(2)$: По определению интеграла Римана-Стилтьеса: 
        \[\forall \epsilon>0\ \exists\ \delta>0,\ d(T)<\delta: |\sigma_s(f\circ \phi,T,\eta)-I|<\frac{\epsilon}{2}\] 
        то есть 
        \[\left|\sum\limits_{i=1}^{n}f(\phi(\eta_i))\cdot(s(t_i)-s(t_{i-1}))-I\right|<\frac{\epsilon}{2}\]
        Обозначим $\xi_i=\phi(\eta_i)$. Тогда
        \[\sigma(f,\Lambda(T),H)=\sum\limits_{i=1}^{n}f(\xi_i)\cdot |\Delta_i|\]
        Пусть $M=\sup\limits_{\gamma}|f|$
        \[\forall \epsilon>0\ \exists\ \theta_1>0,\ d(\Lambda(T))<\theta_1:\ |\gamma|-\Lambda(T)<\frac{\epsilon}{2M}\]
        \begin{multline*}
            |\sigma_s-\sigma|=\left|\sum\limits_{i=1}^{n}f(\xi_i)\cdot(s(t_i)-s(t_{i-1})-|\Delta_i|)\right|\leq\\
            \leq M\cdot \sum\limits_{i=1}^{n}(s(t_i)-s(t_{i-1})-|\Delta_i|)\leq M(|\gamma|-\Lambda(T))<M\cdot\frac{\epsilon}{2M}=\frac{\epsilon}{2}
        \end{multline*}
        Заметим, что отображение $\phi^{-1}$ непрерывно на компакте, значит, оно равномерно непрерывно. Таким образом, если $C,D\in \gamma$ такие, что $\rho(C,D)<\theta_2$, то $|\phi^{-1}(C)-\phi^{-1}(D)|<\delta$. Выберем $\theta=\min\{\theta_1,\theta_2\}$, тогда при $d(\Lambda(T))<\theta$:
        \[|\sigma-I|<\epsilon\]
        \item $(2)\Rightarrow(1)$: Аналогично.
    \end{enumerate}
\end{proof}
%тут сформулировали свойства 0-3